% Generated by tools/edn_latex.py from reports/edn.md.
% Do not edit: change reports/edn.md and run `make edn`.
\ednentry{
  id = {EDN-34},
  title = {The derived raw Parquet keeps its PII columns and is pushed like any other stage output},
  date = {2026-09-30},
  milestone = {M2: Reproducibility},
  topic = {Data Versioning, Personal Data},
  participants = {@lukas2510. Open to revisit at the next sprint review, like EDN-25.},
  decision = {\texttt{data/\allowbreak{}raw/\allowbreak{}listings.\allowbreak{}parquet}, the output of the \texttt{download} stage, keeps the seven PII columns of the published file and is cached and pushed to the DagsHub remote like any other stage output. No \texttt{push: false}, and no PII removal before preprocessing. Our public remote therefore holds the same personal data twice, once as the CSV of EDN-25 and once as this Parquet.},
  rationale = {\emph{Alternatives considered:}
\begin{itemize}[leftmargin=*, topsep=2pt]
\item \textbf{Option A (chosen): push it, which is DVC's default for a stage output.}
\newline Pros: nothing to build. The marginal exposure is close to nothing, because the identical data already sits on the same public remote as a CSV (EDN-25) and on Zenodo under the same licence, so anyone who wants those columns already has them. \texttt{dvc pull} stays sufficient for a complete working state, and the pipeline's first artefact is byte-identical for everyone because it is pulled rather than regenerated.
\newline Cons: we hold a second artefact containing the same personal data under our own name. A deletion request, a licence question or a scope change then touches two artefacts instead of one. EDN-25's argument was that \emph{the identical file} is already public, and a format we produced ourselves is not literally that file.
\item \textbf{Option B: \texttt{push: false} on the output.} DVC keeps it in the local cache only and each person regenerates it from the CSV.
\newline Pros: no second copy leaves the machine.
\newline Cons: it moves the pipeline's first artefact from pulled to locally produced. Parquet writing is not guaranteed byte-stable across pyarrow versions, so a teammate on a different version would see every downstream stage rerun, which is exactly the reproducibility property NFR-06 asks for and that the skeleton was just measured to have.
\item \textbf{Option C: drop the PII in \texttt{download} and build \texttt{seller\_\allowbreak{}group\_\allowbreak{}id} there.}
\newline Pros: the Parquet every stage reads would carry no PII at all, so NFR-08 would hold from the first derived artefact rather than from the interim frame, and the file would still be pushed, so reproducibility would be unaffected.
\newline Cons: the raw layer stops being a faithful copy of the published file, which is what \texttt{RAW\_\allowbreak{}SCHEMA} and the raw expectation suite of \#25 exist to describe; and the group-key hashing moves out of \#34's first step into \#33, so a ticket boundary shifts for a gain that does not change what is publicly reachable.
\item \textbf{Option D: no Parquet at all, every stage reads the CSV.}
\newline Cons: each run parses 548.6 MB of CSV and re-infers its dtypes, which is the cost the stage exists to remove.
\end{itemize}
\par
\emph{Why:} The question is marginal exposure, not storage: measured on the real file the CSV is 548.6 MB, the Parquet of all 75 columns is 215.3 MB, and the seven PII columns are 3.1 MB of it. Against a copy of the same data that is already public on our own remote and on Zenodo, none of the alternatives changes what a reader can reach; they only change what we have to build and what we risk. B risks the reproducibility we have just demonstrated, and C buys the same nothing in exposure terms while making the raw layer no longer raw. So the honest reading is that A costs the least and hides the least.
\par
What A does cost is recorded rather than argued away: we become the publisher of that personal data in two artefacts instead of one, and EDN-25's ``the identical file is already public'' does not literally cover a format we produced ourselves. The entry exists so that a later decision to restrict the remote, or a request to remove the data, finds both artefacts named in one place.
\par
NFR-08 is unaffected either way. It forbids PII in the processed data, the prediction log, the comparables and the model artefacts; the raw layer is none of those, and preprocessing removes the columns structurally, because \texttt{Schema.\allowbreak{}conform} selects only the columns the interim contract names.},
  ai = {Information seeking, Alternative generation, Alternative assessment, Recommendation},
  response = {Accepted with modifications},
  assessment = {AI found the question in the first place, while reviewing the pipeline skeleton: the \texttt{download} stage's output had inherited EDN-25's acceptance without anyone noticing that EDN-25 had only ever weighed the CSV. Its first write-up estimated the Parquet at ``roughly another 550 MB''; asked to explain the trade-off, it measured instead and reported 215.3 MB, of which the PII columns are 3.1 MB, which moved the argument off storage entirely. It also found, while explaining, a reproducibility flaw in its own earlier suggestion of \texttt{push: false} and withdrew it. It recommended C as the clean option and A as defensible; Lukas chose A as the cheapest of the two it stood behind. The modification is that the cost of A is written down here rather than treated as settled by EDN-25.},
  aievidence = {Claude Code session on 2026-09-29 and 2026-09-30: the question was raised in the review of PR \#48, the sizes were measured against the Zenodo file, and the four options were laid out with the reproducibility catch in B; Lukas replied ``mache die einfachste variante die du recommendest''.},
  evidence = {\hyperref[EDN-25]{EDN-25}; \hyperref[EDN-20]{EDN-20}; \href{https://github.com/mlops-2627q1-mds-upc/MLOPS_Recommenditos/issues/33}{issue \#33}; \href{https://github.com/mlops-2627q1-mds-upc/MLOPS_Recommenditos/pull/48}{PR \#48}.},
}
